\documentclass{beamer}
\usepackage{amsmath}
\usepackage{amssymb}

% Choose a theme
\usetheme{Madrid}
\usecolortheme{default}
\usefonttheme{serif}

% --- CUSTOM FOOTER ---
% Redefine the footline to remove author/institute and keep title/page number
\makeatletter
\setbeamertemplate{footline}
{
  \leavevmode%
  \hbox{%
  \begin{beamercolorbox}[wd=.5\paperwidth,ht=2.25ex,dp=1ex,center]{title in head/foot}%
    \usebeamerfont{title in head/foot}\insertshorttitle
  \end{beamercolorbox}%
  \begin{beamercolorbox}[wd=.5\paperwidth,ht=2.25ex,dp=1ex,right]{date in head/foot}%
    \usebeamerfont{date in head/foot}\insertframenumber{} / \inserttotalframenumber\hspace*{2ex}
  \end{beamercolorbox}}%
  \vskip0pt%
}
\makeatother
% --- END CUSTOM FOOTER ---

\title[Lecture 23: Mode splitting and coupling]{Lecture 23: Mode splitting and coupling}
\author{David Al-Attar}
\institute{Michaelmas term 2024}
\date{}

\begin{document}

% --- INTRODUCTION ---

\begin{frame}
    \titlepage
\end{frame}

\begin{frame}{Outline and Motivation}
    In this final lecture, we take a closer look at the Earth's free oscillations.
    \begin{itemize}
        \item We will first classify the normal modes in a perfectly spherically symmetric Earth, where they are highly \textbf{degenerate}.
        \pause
        \item We will then use first-order perturbation theory to show how this degeneracy is broken or "split" by the effects of the Earth's \textbf{rotation} and \textbf{lateral heterogeneity}.
        \pause
        \item A key result, the \textbf{diagonal sum rule}, will show us how to use observations of split modes to learn about the Earth's average radial structure.
        \pause
        \item This provides the primary method for constraining the Earth's density profile.
    \end{itemize}
\end{frame}

% --- MODE CLASSIFICATION ---

\begin{frame}{Modes in a Spherical Earth: Two Families}
    In a spherically symmetric Earth, the normal mode eigenfunctions separate into two distinct families based on their motion.
    \begin{columns}
        \column{0.5\textwidth}
        \begin{block}{Toroidal Modes}
            \begin{itemize}
                \item Purely tangential motion.
                \item Like twisting a rubber ball.
                \item No radial motion, no volume change.
                \item No perturbation to the gravitational field.
                \item Cannot exist in fluid regions (like the outer core).
            \end{itemize}
        \end{block}
        \column{0.5\textwidth}
        \begin{block}{Spheroidal Modes}
            \begin{itemize}
                \item Both radial and tangential motion.
                \item Like breathing or squeezing a ball.
                \item Involve compression and dilatation.
                \item Perturb the gravitational field.
                \item Can exist in both solid and fluid regions.
            \end{itemize}
        \end{block}
    \end{columns}
\end{frame}

\begin{frame}{Mode Classification and Degeneracy}
    Each normal mode is identified by a set of integer indices:
    \begin{itemize}
        \item $ n $: The \textbf{overtone number} ($n=0, 1, 2, \dots$), related to the number of radial nodes.
        \item $ l $: The \textbf{angular degree} ($l=0, 1, 2, \dots$), related to the number of nodal lines on the surface.
        \item $ m $: The \textbf{azimuthal order} ($-l \le m \le l$), related to the pattern of motion around the rotation axis.
    \end{itemize}
    \pause
    \begin{alertblock}{Degeneracy}
        In a perfectly spherical, non-rotating Earth, the frequency of a mode depends on $n$ and $l$, but \textbf{not on $m$}.
        \begin{itemize}
            \item For a given $n$ and $l$, there are $2l+1$ modes (one for each $m$) that all share the exact same frequency.
            \item This is a $(2l+1)$-fold \textbf{degeneracy}.
            \item The set of these $2l+1$ modes is called a \textbf{multiplet} (e.g., ${}_0S_2$). The individual modes within it are called \textbf{singlets}.
        \end{itemize}
    \end{alertblock}
\end{frame}

% --- MODE SPLITTING ---

\begin{frame}{Breaking the Degeneracy: Perturbation Theory}
    The Earth is not perfectly spherical or non-rotating. These aspherical perturbations break the degeneracy.
    \begin{itemize}
        \item This is analogous to the \textbf{Zeeman effect} in quantum mechanics, where a magnetic field splits the degenerate energy levels of an atom.
        \pause
        \item We treat rotation and lateral heterogeneity as small perturbations to the spherical reference model.
        \pause
        \item We use \textbf{degenerate perturbation theory} to find the first-order corrections, $\delta\omega$, to the eigenfrequencies.
    \end{itemize}
    \pause
    \begin{block}{Result}
        The frequency shifts for the $2l+1$ singlets within a multiplet are found by solving a $(2l+1) \times (2l+1)$ matrix eigenvalue problem.
    \end{block}
\end{frame}

\begin{frame}{The Splitting Matrix}
    The perturbation theory leads to the following eigenvalue problem for the frequency shifts $\delta\omega$:
    \begin{alertblock}{The Eigenvalue Problem for Splitting}
        $$ \mathbf{H}' \mathbf{a} = \delta\omega \mathbf{a} $$
    \end{alertblock}
    \begin{itemize}
        \item $\mathbf{H}'$ is a $(2l+1) \times (2l+1)$ Hermitian matrix whose elements measure the interaction of the perturbation with the modes.
        \pause
        \item An element $H'_{m'm}$ represents how the perturbation couples singlet $m$ with singlet $m'$.
        \pause
        \item The $2l+1$ eigenvalues of this matrix are the frequency shifts $\delta\omega$ for each of the singlets. The eigenvectors $\mathbf{a}$ describe how the original degenerate eigenfunctions mix.
        \pause
        \item In general, the eigenvalues will be distinct, so the multiplet "splits" into $2l+1$ observable peaks in the spectrum.
    \end{itemize}
\end{frame}

\begin{frame}{Selection Rules}
    Many of the elements in the splitting matrix are zero due to symmetry. The conditions for an element to be non-zero are called \textbf{selection rules}.
    \begin{itemize}
        \item These rules arise from the properties of integrals of spherical harmonics.
        \pause
        \item \textbf{Example 1:} A structural heterogeneity with angular degree $s$ can only cause splitting within a multiplet of degree $l$ if $s \le 2l$ and $s$ is \textbf{even}. Odd-degree structure does not cause first-order splitting.
        \pause
        \item \textbf{Example 2:} Rotation acts like a perturbation of degree $s=2$ (due to the centrifugal bulge) and $s=1$ (Coriolis force, in a more advanced treatment).
    \end{itemize}
    \pause
    \begin{exampleblock}{Analogy to Quantum Mechanics}
        These rules are mathematically identical to the selection rules governing transitions and interactions in quantum systems involving angular momentum.
    \end{exampleblock}
\end{frame}

% --- DENSITY STRUCTURE ---

\begin{frame}{The Diagonal Sum Rule}
    A key result from splitting theory is the \textbf{diagonal sum rule} (Gilbert, 1971).
    \begin{block}{The Rule}
        The trace of the splitting matrix $\mathbf{H}'$ is zero. Since the trace of a matrix is the sum of its eigenvalues, this means:
        $$ \sum_{m=-l}^{l} \delta\omega_m = 0 $$
        The sum of the frequency shifts of all singlets within a multiplet is zero.
    \end{block}
    \pause
    \begin{alertblock}{Practical Implication}
        This gives us a powerful tool. We can observe the frequencies of the split singlets. By taking their arithmetic average, we can recover the frequency of the unsplit multiplet, which is the frequency the Earth would have if it were perfectly spherical.
        $$ \omega_{unsplit} \approx \frac{1}{2l+1} \sum_{m=-l}^{l} \omega_{obs, m} $$
    \end{alertblock}
\end{frame}

\begin{frame}{Constraining the Earth's Density}
    \begin{itemize}
        \item The diagonal sum rule allowed seismologists to compile a large dataset of unsplit multiplet frequencies from observed spectra of large earthquakes.
        \pause
        \item Travel times of seismic waves are primarily sensitive to wave \textit{speed}.
        \pause
        \item Normal mode frequencies, because they depend on self-gravitation, are sensitive to both wave \textbf{speed} and \textbf{density}.
    \end{itemize}
    \pause
    \begin{alertblock}{A Key Result}
        By jointly inverting travel-time data and normal mode frequency data, it became possible to independently solve for the Earth's radial velocity and density profiles. This is how we know the density of the core and mantle.
    \end{alertblock}
\end{frame}

\begin{frame}{Lateral Density Variations and Mode Coupling}
    A major goal of modern seismology is to map 3D variations in density, which is crucial for understanding mantle dynamics.
    \begin{itemize}
        \item For example, are the Large Low Shear Velocity Provinces (LLSVPs) at the base of the mantle hot and buoyant, or are they compositionally distinct and dense?
        \pause
        \item Simple perturbation theory is often not accurate enough for this, because it ignores the coupling between different multiplets.
        \pause
        \item Modern research uses \textbf{mode coupling theory}, which solves a much larger eigenvalue problem that includes the interactions between nearby multiplets.
        \pause
        \item This is computationally very expensive, and constraining 3D density variations remains a challenging frontier of seismology.
    \end{itemize}
\end{frame}

% --- SUMMARY ---

\begin{frame}{Summary: What You Need to Know}
    \begin{itemize}
        \item In a spherical Earth, normal modes are classified as \textbf{toroidal} or \textbf{spheroidal} and are highly \textbf{degenerate}.
        \pause
        \item Rotation and lateral heterogeneity break this degeneracy, \textbf{splitting} a multiplet into $2l+1$ observable singlets. This can be calculated using degenerate perturbation theory.
        \pause
        \item \textbf{Selection rules}, based on spherical harmonic properties, determine which perturbations can cause splitting and coupling.
        \pause
        \item The \textbf{diagonal sum rule} states that the average frequency of a split multiplet is equal to the unsplit frequency of the reference spherical Earth.
        \pause
        \item This rule was crucial for using observed mode frequencies to constrain the Earth's radial \textbf{density} structure.
    \end{itemize}
\end{frame}

\end{document}
